\chapter{Introduction}
\label{sec:introduction}

In early 2025, a curious pattern emerged across the multimodal AI landscape. Models like GPT-4o and Gemini-1.5 could describe images in exquisite detail, answer complex questions about charts and documents, and even reason about abstract visual puzzles, yet they struggled with questions that a child could answer intuitively. \textit{``Which object is closer to the camera?''} \textit{``If you turned around, what would you see?''} \textit{``How far apart are the chair and the desk?''} These are not exotic tasks; they are the bread and butter of how humans navigate the physical world.

This observation motivated the research that this report documents. The central thesis is simple but important. Spatial intelligence is both critically important and underserved by current evaluation and training paradigms. The work proceeds in three phases. First, we ask, \textit{how bad is the problem, really?} To answer this rigorously requires a holistic evaluation framework that goes beyond individual benchmarks. Second, armed with a precise diagnosis of what models get wrong, we ask, \textit{can we fix it through data?} Specifically, can curating and scaling spatial training data systematically improve these capabilities? Third, and most ambitiously, \textit{does any of this matter for agents that must act in the real world?} Static benchmarks test perception and reasoning in isolation, but embodied agents must also plan, execute, and recover from failure in interactive 3D environments.

Three questions---diagnosis, treatment, and validation---structure the report.
